\section{Introduction}\label{sec:intro}

AlphaFold2 predicts the three-dimensional structure of a protein from its amino-acid sequence~\cite{jumper2021alphafold}.
Its public implementation is written in JAX~\cite{alphafoldrepo}. JAX compiles the model with the accelerated linear
algebra (XLA) compiler for central processing units (CPUs), graphics processing units (GPUs), and tensor processing
units (TPUs)~\cite{frostig2018jax}. A TPU slice connects several TPU chips, and JAX treats each chip as a separate
device. A JAX program on a slice does not spread its work over every chip by itself. The program has to map its work
to the chips, for example with \texttt{jax.pmap}~\cite{jaxpmap}. At list price, a slice costs the same whether one
chip works or all of them do (Section~\ref{sec:res-cost}). So a user who runs AlphaFold2 on a slice without mapping
the work can pay for chips that do nothing.

Our goal is to answer two questions for a user of AlphaFold2 on a TPU slice. First, how much of the slice does
AlphaFold2 inference use by default? Second, which way of mapping the work in JAX uses the rest, and how does
throughput grow with the number of chips?

This study began as a project in a high performance computing course at Stanford University. TPUs were a main topic of
the course. After the course, the project grew into the study we report here.
We ran AlphaFold2 inference on an eight-chip TPU v5e slice of the course cluster in 12 experiments. % src: recomputed: COUNT.tpu_experiments_with_raw_data
We compared four ways of mapping the
work. They are the default, batching with \texttt{jax.vmap}, multi-query \texttt{jax.pmap}, and automatic
partitioning. The model used random weights, so the study concerns execution and not the quality of the predicted
structures.

Work on AlphaFold2 performance has mostly targeted GPUs. OpenFold reimplements AlphaFold2 in
PyTorch~\cite{ahdritz2024openfold}. FastFold runs AlphaFold2 on GPU clusters~\cite{cheng2024fastfold}. ParaFold splits
the pipeline between CPU nodes and GPUs~\cite{zhong2022parafold}. It sorts its input sequences by length, so that JAX
compiles the model once for many of them~\cite{zhong2022parafold}. ColabFold speeds up batch predictions about 90-fold, in part by avoiding recompilation~\cite{colabfold2022}. % src: colabfold2022 ("speeds up batch predictions by ~90-fold by avoiding recompilation and adding an early stop criterion"), paper-v3/SOURCE_NOTES.md section 1
APACE runs AlphaFold2 on 200 NVIDIA A100 GPUs~\cite{apace2024}. % src: apace2024 abstract ("Using up to 300 ensembles, distributed across 200 NVIDIA A100 GPUs"), paper-v3/SOURCE_NOTES.md section 1
ManyFold trains a protein folding model on TPUs~\cite{villegasmorcillo2023manyfold}. JAXBench includes one AlphaFold2
operator among its TPU kernels~\cite{jaxbench2026}. A concurrent working paper measures ColabFold and AlphaFold2 inference on one
Cloud TPU v5e or v6e chip~\cite{forsyth2026timetoscience}. We found no published measurement of AlphaFold2 inference
across the chips of a multi-chip TPU slice other than versions 1 and 2 of this paper.
On 23 November 2026, we searched OpenAlex, DBLP, arXiv, IEEE Xplore, Crossref, Europe PMC, the Google Cloud blog, and the GitHub issues of AlphaFold2 and ColabFold. % src: paper-v3/LITERATURE_SEARCH.md ("23 November 2026")

In this paper, a query is one input sequence that the model predicts once.
The study also has one baseline configuration (Table~\ref{tab:experiments}).
The baseline configuration ran on the TPU, on an NVIDIA T4 GPU, and on a Google Colab CPU runtime.

By default, AlphaFold2 ran on one chip of the slice, and the other seven chips held no data. Batching up to eight
queries on one chip with \texttt{jax.vmap} did not raise throughput. Mapping independent queries to the chips with
\texttt{jax.pmap} ran them in parallel on all eight chips. On a grid of chip counts and sequence lengths, eight chips gave 6.53 to 7.91 times the throughput of one chip. % src: results/sweep/scaling_law_data.json -> grid; recomputed: V2.speedup_8over1_min, V2.speedup_8over1_max
JAX's automatic partitioning did not shorten the call for one query. We also built an external ensemble, which runs
differently featurized copies of one query outside the model. It ran on all eight chips. A compilation cache shared by
two processes shortened parameter initialization and the first call of the second process.

We report the comparison with the T4 and the CPU runtime as context, not as a result. The three platforms ran in
different sessions and with different copies of the benchmark script. On the T4 and on the CPU runtime, the times
changed between sessions (Section~\ref{sec:res-platforms}). The results on the slice come from one platform and one period.

% The sentence never ends a column alone. \keepwithlist (macros.tex) moves it to the next column
% together with the first item of the list when they do not fit. The text and the indent do not change.
\keepwithlist
This paper makes three contributions.
\begin{itemize}
\item We measure how AlphaFold2 inference uses an eight-chip TPU v5e slice under the four ways of mapping the work
described above (Sections~\ref{sec:res-default} to~\ref{sec:res-auto}).
\item We give the throughput on a grid of chip counts and sequence lengths when each chip runs one query through
\texttt{jax.pmap}. We fit a descriptive power law to this grid and give the cost of the queries at list price
(Sections~\ref{sec:res-pmap} and~\ref{sec:res-cost}).
\item We list the script, the input, the timed region, and the number of runs of every experiment. We also state
which records were not kept (Section~\ref{sec:method}).
\end{itemize}

This paper is version 3 of arXiv:2609.34818.
\ifappendix Appendix~\ref{app:changes} lists the claims of the earlier versions that changed.\fi
